\oldpage[I]{80}
\begin{env}[1.1.1]
\label{I.1.1.1-fr}
\emph{Notations}~: Soient $A$ un anneau (commutatif), $M$ un $A$-module. Dans
ce chapitre et les suivants, nous utiliserons constamment les notations
suivantes~:

\smallskip
\noindent
$\operatorname{Spec}(A)={}$\emph{ensemble des idéaux premiers} de $A$, appelé
aussi \emph{spectre premier} de $A$~; pour un
$x\in X=\operatorname{Spec}(A)$, il sera souvent commode d'écrire
$\mathfrak{j}_x$ au lieu de $x$. Pour que $\operatorname{Spec}(A)$ soit
\emph{vide}, il faut et il suffit que l'anneau $A$ soit réduit à $0$.

\smallskip
\noindent
$A_x=A_{\mathfrak{j}_x}={}$\emph{anneau (local) des fractions $S^{-1}A$}, où
$S=A-\mathfrak{j}_x$.

\smallskip
\noindent
$\mathfrak{m}_x=\mathfrak{j}_xA_{\mathfrak{j}_x}={}$\emph{idéal maximal de
$A_x$}.

\smallskip
\noindent
$k(x)=A_x/\mathfrak{m}_x={}$\emph{corps résiduel de $A_x$}, isomorphe
canoniquement au corps des fractions de l'anneau intègre
$A/\mathfrak{j}_x$, auquel on l'identifie.

\smallskip
\noindent
$f(x)={}$\emph{classe de $f$ mod.~$\mathfrak{j}_x$}, dans
$A/\mathfrak{j}_x\subset k(x)$, pour $f\in A$ et $x\in X$. On dit encore que
$f(x)$ est la \emph{valeur} de $f$ au point $x\in\operatorname{Spec}(A)$~;
les relations $f(x)=0$ et $f\in\mathfrak{j}_x$ sont \emph{équivalentes}.

\smallskip
\noindent
$M_x=M\otimes_AA_x={}$\emph{module des fractions à dénominateurs dans
$A-\mathfrak{j}_x$}.

\smallskip
\noindent
$\mathfrak{r}(E)={}$\emph{racine de l'idéal de $A$ engendré par une partie
$E$ de $A$}.

\smallskip
\noindent
$V(E)={}$\emph{ensemble des $x\in X$ tels que $E\subset\mathfrak{j}_x$}
(ou encore ensemble des $x\in X$ tels que $f(x)=0$ pour tout $f\in E$), pour
$E\subset A$. On a donc
\begin{equation}
  \mathfrak{r}(E)=\bigcap_{x\in V(E)}\mathfrak{j}_x.
  \tag{1.1.1.1}\label{I.1.1.1.1-fr}
\end{equation}

\smallskip
\noindent
$V(f)=V(\{f\})$ pour $f\in A$.

\smallskip
\noindent
$D(f)=X-V(f)={}$\emph{ensemble des $x\in X$ où $f(x)\neq0$}.
\end{env}

\begin{proposition}[1.1.2]
\label{I.1.1.2-fr}
On a les propriétés suivantes~:
\begin{enumerate}
  \item[\rm(i)] $V(0)=X$, $V(1)=\varnothing$.
  \item[\rm(ii)] La relation $E\subset E'$ entraîne $V(E)\supset V(E')$.
  \item[\rm(iii)] Pour toute famille $(E_\lambda)$ de parties de $A$,
  \[
    V\left(\bigcup_\lambda E_\lambda\right)
      =V\left(\sum_\lambda E_\lambda\right)
      =\bigcap_\lambda V(E_\lambda).
  \]
  \item[\rm(iv)] $V(EE')=V(E)\cup V(E')$.
  \item[\rm(v)] $V(E)=V(\mathfrak{r}(E))$.
\end{enumerate}
\end{proposition}

Les propriétés (i), (ii), (iii) sont triviales, et (v) résulte de (ii) et de
la formule (\hyperref[I.1.1.1.1-fr]{1.1.1.1}). Il est évident que
$V(EE')\supset V(E)\cup V(E')$~; inversement, si $x\notin V(E)$ et
$x\notin V(E')$, il existe $f\in E$ et $f'\in E'$ tels que $f(x)\neq0$ et
$f'(x)\neq0$ dans $k(x)$, d'où $f(x)f'(x)\neq0$, autrement dit
$x\notin V(EE')$, ce qui prouve (iv).

La prop. (\hyperref[I.1.1.2-fr]{1.1.2}) montre entre autres que les ensembles
de la forme $V(E)$ (où $E$ parcourt l'ensemble des parties de $A$) sont les
\emph{ensembles fermés} d'une topologie sur $X$, que nous appellerons la
\emph{topologie spectrale}\footnote{L'introduction de cette topologie en
géométrie algébrique est due à Zariski. Aussi est-elle souvent appelée la
<<~topologie de Zariski~>> de $X$.}~; sauf mention expresse du contraire, on
supposera toujours $X=\operatorname{Spec}(A)$ muni de la topologie spectrale.

